\documentclass{beamer}

\input{MathMacros}

\usepackage{beamerthemesplit}
\usepackage{enumerate}
\usepackage{amssymb}
\usepackage[all,cmtip]{xy}
\usepackage[mathscr]{eucal}
\usepackage[natbib=true,style=authoryear,backend=bibtex,useprefix=true]{biblatex}
\addbibresource{reading.bib}
\usepackage{graphicx,color}

\title{From Scaffolding to Governance: The Future of AI Work}
\author{Reuben Brasher}
\date{\today}

\begin{document}

\frame{\titlepage}

\section[Outline]{}
\frame{\tableofcontents}

\section{Lecture 9}

\frame{
\frametitle{From Tools to Systems}
AI development progresses from local scaffolding to organizational governance.
\begin{itemize}
\item Prompting introduces initial constraints.
\item Workflows structure interactions.
\item Systems coordinate components.
\item Governance formalizes constraints.
\item Usage evolves toward structure.
\item Control scales institutionally.
\end{itemize}
}

\frame{
\frametitle{A Single Trajectory}
The series describes a progression of increasing constraint.
\begin{itemize}
\item Phase one focused on cognition.
\item Phase two introduced workflows.
\item Phase three revealed failures.
\item Phase four introduced governance.
\item Each phase added control.
\item Structure increased progressively.
\end{itemize}
}

\frame{
\frametitle{Core Synthesis}
Constraint is required to stabilize inherently unconstrained systems.
\begin{itemize}
\item LLMs produce unconstrained outputs.
\item Structure limits output space.
\item Constraints improve reliability.
\item Lack of constraints causes failure.
\item Control must be imposed.
\item Stability emerges from layering.
\end{itemize}
}

\frame{
\frametitle{Scaffolding Defined}
Scaffolding enables reliable interaction with models.
\begin{itemize}
\item It guides model behavior.
\item It reduces ambiguity.
\item It aligns representations.
\item It enables refinement.
\item It supports validation.
\item It acts as control.
\end{itemize}
}

\frame{
\frametitle{Forms of Scaffolding}
Scaffolding appears in multiple forms.
\begin{itemize}
\item Priming establishes context.
\item Schema defines structure.
\item Decomposition breaks tasks.
\item Testing verifies correctness.
\item Role separation isolates work.
\item Each form constrains failure.
\end{itemize}
}

\frame{
\frametitle{Reliability Through Scaffolding}
Reliability depends on structured scaffolding.
\begin{itemize}
\item Unstructured prompts are unstable.
\item Workflows improve consistency.
\item Validation prevents errors.
\item Decomposition reduces complexity.
\item Role separation limits propagation.
\item Structure stabilizes outputs.
\end{itemize}
}

\frame{
\frametitle{Critical Shift}
Scaffolding is the system rather than a workaround.
\begin{itemize}
\item It defines model usage.
\item It determines output quality.
\item It replaces internal reasoning.
\item It creates reproducibility.
\item It enables collaboration.
\item It transforms tools into systems.
\end{itemize}
}

\frame{
\frametitle{From Prompts to Workflows}
Prompts evolve into structured workflows.
\begin{itemize}
\item Prompts become multi-step processes.
\item Steps are sequenced.
\item Outputs feed forward.
\item Validation is inserted.
\item Roles are assigned.
\item Workflows become repeatable.
\end{itemize}
}

\frame{
\frametitle{From Workflows to Pipelines}
Workflows scale into pipelines.
\begin{itemize}
\item Processes become standardized.
\item Steps are automated.
\item Data flows are explicit.
\item Dependencies are managed.
\item Outputs are structured.
\item Pipelines enable scale.
\end{itemize}
}

\frame{
\frametitle{From Pipelines to Systems}
Pipelines evolve into integrated systems.
\begin{itemize}
\item Multiple pipelines interact.
\item Data flows across components.
\item Systems integrate infrastructure.
\item Monitoring is introduced.
\item Control is formalized.
\item Systems operate continuously.
\end{itemize}
}

\frame{
\frametitle{Key Insight}
Prompting is an early form of system design.
\begin{itemize}
\item Prompts define inputs.
\item Workflows define logic.
\item Validation defines correctness.
\item Roles define architecture.
\item Iteration defines development.
\item Design emerges from interaction.
\end{itemize}
}

\frame{
\frametitle{Transition to Governance}
Systems that process data must be governed.
\begin{itemize}
\item Data flows require control.
\item Outputs require accountability.
\item Processes require documentation.
\item Systems must meet regulation.
\item Governance ensures consistency.
\item Control scales organizationally.
\end{itemize}
}

\frame{
\frametitle{Scaffolding to Governance}
Local scaffolding maps to governance structures.
\begin{itemize}
\item Prompts become policies.
\item Schema becomes standards.
\item Testing becomes auditing.
\item Roles become organizational roles.
\item Validation becomes compliance.
\item Workflows become governed processes.
\end{itemize}
}

\frame{
\frametitle{Governance as Scaffolding}
Governance is system-level scaffolding.
\begin{itemize}
\item Defines system behavior.
\item Enforces constraints.
\item Standardizes practices.
\item Ensures compliance.
\item Enables accountability.
\item Sustains reliability.
\end{itemize}
}

\frame{
\frametitle{Reinterpreting Agentic AI}
Agentic AI is structured and governed workflow execution.
\begin{itemize}
\item Agents execute processes.
\item Memory provides state.
\item Constraints define actions.
\item Outputs are pipeline artifacts.
\item Behavior must be controlled.
\item Autonomy is bounded.
\end{itemize}
}

\frame{
\frametitle{Core Definition of an Agent}
An agent is a constrained pipeline with memory.
\begin{itemize}
\item It receives defined inputs.
\item It processes structured steps.
\item It maintains state.
\item It produces constrained outputs.
\item It operates under governance.
\item It is not autonomous.
\end{itemize}
}

\frame{
\frametitle{Hype vs Reality}
Agent narratives obscure system realities.
\begin{itemize}
\item Hype emphasizes autonomy.
\item Reality emphasizes control.
\item Hype suggests independence.
\item Reality requires supervision.
\item Hype ignores governance.
\item Reality depends on it.
\end{itemize}
}

\frame{
\frametitle{Future Skill Shift}
AI changes required competencies.
\begin{itemize}
\item Less focus on coding.
\item Less focus on basic writing.
\item More focus on structuring problems.
\item More focus on workflow design.
\item More focus on validation.
\item More focus on systems thinking.
\end{itemize}
}

\frame{
\frametitle{New Core Competencies}
Expertise shifts toward design and supervision.
\begin{itemize}
\item Define clear boundaries.
\item Design workflows.
\item Evaluate outputs.
\item Detect errors.
\item Integrate systems.
\item Maintain alignment.
\end{itemize}
}

\frame{
\frametitle{Emerging Roles}
New roles reflect structured interaction.
\begin{itemize}
\item AI architects design systems.
\item Reviewers validate outputs.
\item Governance specialists enforce policy.
\item Engineers integrate systems.
\item Compliance ensures alignment.
\item Roles mirror workflows.
\end{itemize}
}

\frame{
\frametitle{Human–AI Collaboration}
LLMs act as junior collaborators.
\begin{itemize}
\item Models generate outputs.
\item Humans define constraints.
\item Humans validate correctness.
\item Humans make decisions.
\item Collaboration is iterative.
\item Responsibility remains human.
\end{itemize}
}

\frame{
\frametitle{Boundaries of Autonomy}
AI autonomy remains limited.
\begin{itemize}
\item Systems operate within constraints.
\item Outputs require validation.
\item Decisions need oversight.
\item Errors are externally detected.
\item Governance defines limits.
\item Autonomy is constrained.
\end{itemize}
}

\frame{
\frametitle{Organizational Implications}
AI systems must align with enterprise environments.
\begin{itemize}
\item Integration with infrastructure is required.
\item Compliance must be maintained.
\item Security must be ensured.
\item Collaboration is necessary.
\item Governance must be embedded.
\item Systems must align with constraints.
\end{itemize}
}

\frame{
\frametitle{Organizational Risk}
Ignoring governance leads to failure.
\begin{itemize}
\item Compliance violations may occur.
\item Security risks increase.
\item Outputs may be unreliable.
\item Stakeholders reject systems.
\item Rework becomes necessary.
\item Projects fail.
\end{itemize}
}

\frame{
\frametitle{Strategic Positioning}
Value lies in system design rather than models.
\begin{itemize}
\item Models are commoditized.
\item Integration creates differentiation.
\item Control ensures reliability.
\item Governance enables adoption.
\item Systems drive value.
\item Strategy focuses on constraints.
\end{itemize}
}

\frame{
\frametitle{Differentiation from Hype}
Model-centric thinking fails in practice.
\begin{itemize}
\item Performance metrics are insufficient.
\item Ignoring governance limits adoption.
\item Autonomy narratives mislead.
\item Systems thinking adds value.
\item Integration matters more.
\item Reliability outweighs capability.
\end{itemize}
}

\frame{
\frametitle{Final Synthesis}
Scaffolding evolves into governance systems.
\begin{itemize}
\item Starts with prompts.
\item Expands into workflows.
\item Scales into pipelines.
\item Becomes governance.
\item Enables reliability.
\item Defines future systems.
\end{itemize}
}

\frame{
\frametitle{Closing Statement}
The future of AI is governed collaboration.
\begin{itemize}
\item Intelligence is insufficient.
\item Structure enables reliability.
\item Governance ensures safety.
\item Collaboration defines effectiveness.
\item Systems determine outcomes.
\item The future is constrained.
\end{itemize}
}

\begin{frame}[t,allowframebreaks]
\frametitle{References}
\printbibliography
\end{frame}

\end{document}